\section{模型评价与推广}

\subsection{模型的优点}

\begin{enumerate}
\item \textbf{数据驱动、可复算}：四问全部数值结果由真实附件数据计算，代码
随附录提供，任何结论均可复现；关键参数（$\alpha,\beta$）与文献报告值
（Chinchilla 的 0.34/0.28）一致，验证了数据与方法的可靠性；
\item \textbf{广义化框架}：把经典标度律推广到含质量维度，且 $Q=1$ 自动
退化，理论一致性良好；同时保证``质量''通过可解释的配比加权进入模型，
贯通问题一--三；
\item \textbf{严谨的数据口径核验}：诊断出附件 B8 的 $Q$ 方向与 B6/B7
相反、C3 历史前沿含孤立高值等数据问题，并给出规避与对照处理，避免了
半合成数据对结论的污染；
\item \textbf{结构性转移的显式数学刻画}：给出质量激活/成本主导两类转移
指标与 KKT 识别方法，使``质变''判断可计算、可复现；
\item \textbf{预测带不确定性}：前沿预测采用自助置信区间与双情景对比，
不给出伪精度的单点预言。
\end{enumerate}

\subsection{模型的缺点与改进方向}

\begin{enumerate}
\item \textbf{半合成数据的局限}：B 族 $N$--$D$--$Q$ 数据为半合成实验，
B8 更与 B6/B7 口径不一致；改进方向：使用真实多质量档预训练实验校准 $Q$
项参数，并做交叉验证；
\item \textbf{广义标度律的可加性假设}：质量项与规模项的可加分解为近似，
当数据量极小或质量极低时交互效应可能显著；改进方向：引入 $N$--$Q$
交互项并比较拟合；
\item \textbf{质量成本函数参数的外生性}：$g(Q)$ 三种形式的 $\gamma,\lambda$
由题意给定且跨度大，实际成本取决于标注、过滤与合成支出；改进方向：用
真实 token 质量提升工程的成本数据标定；
\item \textbf{前沿预测的线性外推}：$\ln S\sim\ln N+t$ 的线性假设在 24 个月
后可能因算法饱和而失效；改进方向：引入饱和项（如 logit 上界）或滚动
再估计窗口；
\item \textbf{排行榜指标的选择偏差}：C1 的平均分受模型提交意愿、评测版本
影响；改进方向：用 C8 逐任务数据构建去偏的因子得分。
\end{enumerate}

\subsection{模型推广}

本文方法可推广到：多模态模型的数据质量评分与配比优化、RAG/微调阶段的质量
--成本权衡、企业算力采购的预算分配决策，以及``数据工程投入 vs 算力投入''
的一般资源规划问题。跨尺度系数收缩分析亦可用于其他``定量投入--性能''系统
（如推荐系统特征工程、模拟器 fidelity 选择），识别投入边际收益随规模的
衰减规律。